% ---------------------------------------------------------------------------
% Code snippet: network summary helper and primary-threshold selection
% Source files: src/network_utils.py, src/02_build_network.py
% Suggested placement: Section 2.2 (Correlation and network construction).
% Uses the shared 'python' listings style defined in main.tex.
% ---------------------------------------------------------------------------
\begin{lstlisting}[style=python]
def graph_metrics(graph: nx.Graph) -> dict:
    components = sorted((len(c) for c in nx.connected_components(graph)),
                        reverse=True)
    n = graph.number_of_nodes()
    return {
        "nodes": n,
        "edges": graph.number_of_edges(),
        "density": nx.density(graph),
        "avg_degree": 2 * graph.number_of_edges() / n if n else 0.0,
        "components": len(components),
        "giant_size": components[0] if components else 0,
        "giant_fraction": components[0] / n if n else 0.0,
        "isolates": nx.number_of_isolates(graph),
    }


def choose_threshold(corr: np.ndarray) -> tuple[float, list[dict]]:
    """Most stringent r that still keeps a usable giant component."""
    scan = [
        graph_metrics(nx.from_numpy_array(adjacency(corr, r)))
        for r in config.THRESHOLD_GRID          # 0.30, 0.31, ..., 0.75
    ]
    feasible = [
        (r, m) for r, m in zip(config.THRESHOLD_GRID, scan)
        if m["avg_degree"] >= config.MIN_AVG_DEGREE
        and m["giant_fraction"] >= config.MIN_GIANT_FRACTION
    ]
    # (the density fallback for an empty feasible set is omitted here)
    return max(feasible, key=lambda x: x[0])[0], scan
\end{lstlisting}
